\chapter{Using This Book in a Course}
\label{app:teaching}

This book was written to be taught. It can support a semester course in a law school or a statistics or economics department, and parts of it can be used in shorter courses or in continuing education for judges and practitioners. We suggest three paths.

\section{A law school seminar on evidence and probability}

A fifteen week seminar for law students, at the level of an upper year law school course, might proceed as follows. Each week combines a chapter with one or two leading cases, which students read in full, and a set of Notes and Questions on which they write short memoranda.
\begin{enumerate}[label=Week \arabic*.,leftmargin=4.6em]
\item Why the law needs Bayes; decisions, objectives, Type I and Type II errors (Chapters~\ref{ch:why} and~\ref{ch:decisions}).
\item Odds, likelihood ratios and the fallacies; \emph{Collins}, \emph{Clark}, \emph{Adams} (Chapter~\ref{ch:odds}).
\item Simpson, the island problem, Balding and Donnelly (Chapter~\ref{ch:classic}).
\item Bayes factors and the fallacies of classical inference; \emph{Matrixx} (Chapters~\ref{ch:bf} and~\ref{ch:fallacies}).
\item Standards of proof; \emph{Re B}, the gatecrasher (Chapter~\ref{ch:standards}).
\item The law of evidence; \emph{Turnbull}, \emph{Hanson}, \emph{Daubert} (Chapter~\ref{ch:evidence}).
\item Tort, mass torts and damages; \emph{Gregg v Scott}, \emph{Fairchild}, \emph{Wells v Wells} (Chapters~\ref{ch:tort}, \ref{ch:masstorts} and~\ref{ch:damages}).
\item Misrepresentation, contracts and Greensill (Chapters~\ref{ch:misrep}, \ref{ch:contract} and~\ref{ch:greensill}).
\item Stories, graphs and pleadings; the Kercher case, \emph{Twombly} (Chapters~\ref{ch:stories} to~\ref{ch:pleadingscases}).
\item Statisticians in court: Fienberg and Kadane; discrimination (Chapters~\ref{ch:fienberg}, \ref{ch:kadane} and~\ref{ch:discrimination}).
\item Clusters and wrongful convictions (Chapter~\ref{ch:clusters}).
\item Comparative law: England, France, Africa and India (Chapters~\ref{ch:commoncivil} to~\ref{ch:india}).
\item Costs, funding, sentencing and regulation (Chapters~\ref{ch:systemcosts} to~\ref{ch:regulation}).
\item Verification and machines (Chapters~\ref{ch:lean} and~\ref{ch:lindleyai}).
\item Explanation, intuition, reform and the future (Chapters~\ref{ch:xai}, \ref{ch:hume}, \ref{ch:program} and~\ref{ch:future}).
\end{enumerate}
No prior statistics is assumed beyond what Appendix~\ref{app:toolkit} supplies. Students who want more mathematics should work the computational questions using the code of Appendix~\ref{app:code}.

\section{A course for statisticians and economists}

For students with a quantitative background, the emphasis shifts to models and computation: Bayes factors and robustness (Chapters~\ref{ch:bf} and~\ref{ch:fallacies}), graphical models and chain event graphs (Chapters~\ref{ch:classic} and~\ref{ch:graphs}), generative Bayesian computation (Chapter~\ref{ch:pleadings}), event studies and heavy tailed priors (Chapter~\ref{ch:securities}), solvency and the Merton model (Chapter~\ref{ch:bankruptcy}), hierarchical models of legal origins (Chapters~\ref{ch:llsv} and~\ref{ch:africa}), queues and triage (Chapter~\ref{ch:india}), the economics of enforcement and antitrust (Chapters~\ref{ch:becker} and~\ref{ch:antitrust}), and verification (Chapter~\ref{ch:lean}). A project might build a Bayesian network for a reported case, reanalyse a published event study, or implement a quantile network for a simulated dispute.

\section{A short course for judges and practitioners}

A two day course for judges, prosecutors and practitioners can be built from Chapters~\ref{ch:odds}, \ref{ch:classic}, \ref{ch:standards}, \ref{ch:evidence}, \ref{ch:clusters}, \ref{ch:hume} and~\ref{ch:lindleyai}, focusing on the fallacies, the reading of forensic reports, the presentation of likelihood ratios to juries, and the use of AI tools in practice. The six step procedure for weighing testimony at the end of Chapter~\ref{ch:hume} makes a good closing exercise: participants apply it to a case from their own experience.

\section{Assessment}

We have found three forms of assessment effective. Short memoranda answering selected Notes and Questions test understanding of the tools. A case analysis, in which students choose a reported case, identify its propositions, evidence and implicit likelihood ratios, and assess whether the court reasoned coherently, tests application. And a reform proposal, in which students draft a rule, direction or practice note implementing one of the proposals of Chapter~\ref{ch:program}, tests judgment. Appendix~\ref{app:exam} gives a sample examination with notes on strong answers.
